\documentclass{article}
\usepackage{amsmath,amssymb,xurl}
\usepackage[hidelinks]{hyperref}
\setlength{\emergencystretch}{2em}
% Shared commands for notes in docs/paper-gaps/.

\DeclareUnicodeCharacter{2080}{\ifmmode{}_{0}\else\textsubscript{0}\fi}
\DeclareUnicodeCharacter{2081}{\ifmmode{}_{1}\else\textsubscript{1}\fi}
\DeclareUnicodeCharacter{2082}{\ifmmode{}_{2}\else\textsubscript{2}\fi}
\DeclareUnicodeCharacter{2099}{\ifmmode{}_{n}\else\textsubscript{n}\fi}

\newcommand{\Lean}{\textsc{Lean}}
\ExplSyntaxOn
\NewDocumentCommand{\leanid}{m}{
  \ifmmode
    \textsf{\detokenize{#1}}
  \else
    \group_begin:
    \urlstyle{sf}
    \tl_set:Nn \l_tmpa_tl {#1}
    \tl_replace_all:Nnn \l_tmpa_tl {\_} {_}
    \exp_args:NV \path \l_tmpa_tl
    \group_end:
  \fi
}
\ExplSyntaxOff
\providecommand{\tr}{\operatorname{tr}}
\newcommand{\tnleanIssueBase}{https://github.com/LionSR/TNLean/issues/}
\newcommand{\tnleanPrBase}{https://github.com/LionSR/TNLean/pull/}
\newcommand{\tnleanDocBase}{https://sirui-lu.com/TNLean/docs/}
\newcommand{\ghissue}[1]{\href{\tnleanIssueBase#1}{GitHub issue \##1}}
\newcommand{\ghpr}[1]{\href{\tnleanPrBase#1}{PR \##1}}
\newcommand{\doclink}[1]{\href{\tnleanDocBase#1.html}{\path{#1}}}

% Dirac notation and the identity operator, as in blueprint/src/macros/common.tex.
% \providecommand keeps note-local definitions authoritative.
\usepackage{dsfont}
\providecommand{\ket}[1]{|#1\rangle}
\providecommand{\bra}[1]{\langle #1|}
\providecommand{\braket}[2]{\langle #1 | #2 \rangle}
\providecommand{\ketbra}[2]{|#1\rangle\!\langle #2|}
\providecommand{\Id}{\mathds{1}}
\providecommand{\one}{\mathds{1}}
\providecommand{\C}{\mathbb{C}}
\providecommand{\MN}[1]{M_{#1}(\C)}
\providecommand{\MD}{M_D(\C)}

% Verdict marker: \gapnote{<kind>}{<status>}. Kind is the policy
% classification (clarification, local-correction, scope-restriction,
% unfaithful, false-source, open-gap); status is open, wip (elimination
% underway), resolved, or historical. Machine-read by the site index;
% prints a verdict line.
\newcommand{\gapnote}[2]{\par\noindent\textbf{Verdict:} #1 (#2).\par}

\title{The ALVV Merge Tree for a Unique Frustration-Free Ground State}
\author{}
\date{November 17, 2026}
\begin{document}
\maketitle
\gapnote{scope-restriction}{open}

\paragraph{Source.}
Arad, Landau, Vazirani and Vidick,
\cite[Theorem~t:unique, lines 1829--1855]{Arad2016RigorousRG}, consider a
nearest-neighbour chain $H=\sum_ih_i$ with $0\le h_i\le\Id$ that is
frustration free, has a unique ground state $\Gamma$ and a spectral gap
$\gamma$. The theorem asserts that the merge-tree algorithm, run with the
frustration-free improving family of
\cite[Theorem~t:aisff, lines 1580--1609]{Arad2016RigorousRG}, outputs a
matrix product state approximating $\Gamma$, in a running time close to that
of matrix multiplication. Its improvement over
\cite[Theorem~t:alg1, lines 1777--1818]{Arad2016RigorousRG} lies in the
running time: the frustration-free family has bond dimension independent of
the chain length, and no energy estimates are carried along the tree.

\paragraph{What is formalized.}
TNLean has no model of computation, so only the existence content is
formalized. The hypotheses of the source are imposed exactly: positive
contractive nearest-neighbour terms, a nonzero vector $\Gamma$ with
$H\Gamma=0$, every vector of the kernel a multiple of $\Gamma$, and every
other spectral value at least $\gamma$. Such a chain satisfies
Assumption~(DG) with ground energy $0$ and ground-space dimension $1$; the
closed spectral span at energy $0$ is the line $\mathbb C\Gamma$. Three
results use this.
\begin{enumerate}
    \item The physical resampling tree of the degenerate-gap case and its
          refinement: on an event of failure probability at most $(L+1)^{-2}$,
          the refinement of the actual root and its Ritz extraction give a line,
          mutually $\eta/2$-close to $\mathbb C\Gamma$, of matrix product states
          whose bond dimension is polynomial in $L$ and
          $1/\eta$.\footnote{\leanid{TNLean.MPS.AreaLaw.ALVV.exists_unique_ground_state_refined_ritz_approximation}.}
          Its node dimension is
          $\max\bigl(q^{\mathrm{width}},D^2K\cdot2\,(L+2)^2\bigr)$, which grows
          with $L$ through the corrected sample threshold of
          \url{https://sirui-lu.com/QICLean/paper-gaps/alvv16_bootstrap_constants.pdf}.
    \item The products of length $k$ of the frustration-free family of
          Theorem~t:aisff improve viability by $\Delta^k/(1-\delta)^2$, as the
          merge of lines 1662--1677
          requires.\footnote{\leanid{TNLean.MPS.AreaLaw.ALVV.exists_frustrationFree_improving_family_words}.}
          Restated under the hypotheses of Theorem~t:unique and in the
          region-first coordinates in which the tree measures viability, this
          carries no hypothesis beyond those of the
          source.\footnote{\leanid{TNLean.MPS.AreaLaw.ALVV.exists_frustrationFree_improving_family_words_regionTensorCoordinates}.}
    \item The deterministic recursive merge tree run with that family on every
          long interval: with the constant $A$ and the parameters
          $g=\min(\gamma,1)$, $d=\max(q,2)$, $u=1+\log d+\log(1/g)$,
          $\ell=\lceil Au^2/g\rceil$, $\mathcal D=(d\ell^2)^{22\ell}$ of
          Theorem~t:aisff, and
          $s=\lceil4096(1+\log\mathcal D+\log(1/\beta))\rceil$ for any
          $0<\beta<1$, every node carries a $3/200$-viable space for
          $\mathbb C\Gamma$ of dimension at most $\mathcal D^6s^3$, and the
          filter-and-trim iterates of the root with a Ritz extraction give a line
          of matrix product states mutually $\eta/2$-close to
          $\mathbb C\Gamma$.\footnote{\leanid{TNLean.MPS.AreaLaw.ALVV.exists_unique_ground_state_tree_ritz_approximation}.}
          Neither $\mathcal D$ nor the node bound depends on $L$.
\end{enumerate}
The former deviation that the tree used the families of the degenerate-gap
case, with constants depending on $\gamma$ through the degenerate-gap
parameters, is eliminated by item~3. The node bound keeps its form
$\mathcal D^6s^3$; with the frustration-free family, $\mathcal D$ is the
parameter of Theorem~t:aisff, whose base is $d=\max(q,2)$ rather than the
$\max(q,50)$ of the degenerate-gap family, so the bound is not sharper in
form.

\paragraph{Deviation.}
The remaining differences from the source are:
\begin{enumerate}
    \item the length-independent node bound is $\mathcal D^6s^3$, not the
          $\mathcal Ds^2$ of the source, and it is attained by a tree shown to
          exist, not by one found with high probability; the probabilistic tree
          of item~1 has node dimension growing with $L$;
    \item the bond dimension of the output matrix product states grows with
          $L$, through the trimming threshold, which is set by a reference space
          whose Schmidt ranks grow with $L$;
    \item item~3 assumes $L\ge1$ and $q\ge2$: on an empty chain, or for $q=1$,
          the space has dimension one, $H$ has no positive spectral value, and the
          refinement of the root does not apply; this excludes only the trivial
          case in which $\mathbb C\Gamma$ is the whole space;
    \item no running time is asserted.
\end{enumerate}
The conclusions are mathematically correct as stated, and narrower than the
source in the items listed.

\paragraph{Elimination plan.}
The constant per-node bound $\mathcal Ds^2$ requires the source's sample
threshold, which is not established.
A reference space for collective trimming whose Schmidt ranks are bounded
independently of $L$ would make the output bond dimension independent of $L$;
it is not needed for the node bound. A probabilistic version of item~3
requires the joint failure bound along the tree.

\bibliographystyle{unsrt}
\bibliography{references}
\end{document}
